\chapter{Supported operators and placement}
\label{ch:sites}

Throughout this chapter $V$ is a finite set of sites and $X=(X_v)_{v\in V}$ is
a family of finite sets. The system with basis sets $X$ is the Hilbert space
$\mathcal H_X=\bigotimes_{v\in V}\C^{X_v}$, whose standard basis is indexed by
the configurations $\sigma\in\prod_{v\in V}X_v$. A region is a subset
$D\subseteq V$; its configurations are the elements of $X_D=\prod_{v\in D}X_v$.
One local dimension $d$ is the case $X_v=\{0,\dots,d-1\}$ for every $v$, and
site-dependent dimensions $n_v$ are the case $X_v=\{0,\dots,n_v-1\}$. No basis
set is required to be nonempty unless stated: if some $X_v$ is empty, then
$\mathcal H_X=0$ and the statements below hold trivially.

\section{Configurations and cuts}
\label{sec:sites_config}

\begin{definition}[Configurations and the cut along a region]
    \label{def:sites_config}
    \lean{ManyBody.Config, ManyBody.RegionConfig, ManyBody.ComplConfig,
        ManyBody.cutEquiv, ManyBody.cutEquiv_symm_apply_mem,
        ManyBody.cutEquiv_symm_apply_not_mem, ManyBody.cutEquiv_fst_eq_iff,
        ManyBody.cutEquiv_snd_eq_iff}
    \leanok
    A configuration of the system is $\sigma\in X_V=\prod_{v\in V}X_v$. For a
    region $D$, the cut along $D$ is the bijection
    $X_V\to X_D\times X_{V\setminus D}$, $\sigma\mapsto(\sigma|_D,
    \sigma|_{V\setminus D})$. Two configurations have the same part on $D$
    exactly when they agree at every site of $D$, and likewise outside $D$.
\end{definition}

\section{Operators supported on a region}
\label{sec:sites_support}

\begin{definition}[Operators supported on a region]
    \label{def:sites_supported}
    \lean{ManyBody.IsSupportedOn, ManyBody.mem_supportedOperators_iff,
        ManyBody.apply_eq_zero_of_mem_supportedOperators,
        ManyBody.apply_eq_of_mem_supportedOperators}
    \leanok
    \uses{def:sites_config}
    An operator $A$ on $\mathcal H_X$ is supported on a region $D$ when its
    matrix element $A_{\sigma\tau}$ vanishes whenever $\sigma_v\neq\tau_v$ for
    some $v\notin D$, and otherwise depends only on $\sigma|_D$ and $\tau|_D$.
    Let $\mathcal A_D$ be the set of operators supported on $D$. Then
    $A_{\sigma\tau}=A_{\sigma'\tau'}$ for $A\in\mathcal A_D$ whenever
    $\sigma|_D=\sigma'|_D$, $\tau|_D=\tau'|_D$ and, at every $v\notin D$,
    $\sigma_v=\tau_v$ if and only if $\sigma'_v=\tau'_v$.
\end{definition}

\begin{theorem}[Supported operators form a unital $*$-subalgebra]
    \label{thm:sites_supported_subalgebra}
    \lean{ManyBody.supportedOperators}
    \leanok
    \uses{def:sites_supported}
    The set $\mathcal A_D$ is a unital $*$-subalgebra of the matrix algebra of
    $\mathcal H_X$.
\end{theorem}
\begin{proof}\leanok
    Sums, scalar multiples and adjoints preserve both conditions entrywise,
    and the identity satisfies them. For a product,
    $(AB)_{\sigma\tau}=\sum_\rho A_{\sigma\rho}B_{\rho\tau}$; only the $\rho$
    that agree with $\sigma$ outside $D$ contribute, so the sum runs over
    $\rho|_D\in X_D$, and each term depends only on the parts on $D$.
\end{proof}

\begin{lemma}[Monotonicity and the full region]
    \label{lem:sites_supported_mono}
    \lean{ManyBody.supportedOperators_mono, ManyBody.supportedOperators_univ}
    \leanok
    \uses{def:sites_supported, thm:sites_supported_subalgebra}
    If $D\subseteq D'$, then $\mathcal A_D\subseteq\mathcal A_{D'}$, and
    $\mathcal A_V$ is the whole matrix algebra.
\end{lemma}
\begin{proof}\leanok
    A difference outside $D'$ is a difference outside $D$, and two pairs with
    the same pattern of coincidences outside $D'$ that agree on $D'$ have the
    same pattern outside $D$.
\end{proof}

\begin{lemma}[Operators on disjoint regions commute]
    \label{lem:sites_supported_commute}
    \lean{ManyBody.commute_of_mem_supportedOperators_of_disjoint,
        ManyBody.mul_apply_of_mem_supportedOperators_of_disjoint}
    \leanok
    \uses{def:sites_supported, thm:sites_supported_subalgebra, lem:sites_supported_mono}
    Let $D\cap D'=\emptyset$, $A\in\mathcal A_D$ and $B\in\mathcal A_{D'}$.
    Then $(AB)_{\sigma\tau}=A_{\sigma\rho}B_{\rho\tau}$, where $\rho$ equals
    $\tau$ on $D$ and $\sigma$ elsewhere, and $AB=BA$.
\end{lemma}
\begin{proof}\leanok
    In $(AB)_{\sigma\tau}=\sum_{\rho}A_{\sigma\rho}B_{\rho\tau}$ the factor
    $A_{\sigma\rho}$ vanishes unless $\rho=\sigma$ outside $D$, and
    $B_{\rho\tau}$ vanishes unless $\rho=\tau$ on $D$. If $\sigma$ and $\tau$
    differ outside $D\cup D'$, both products lie in $\mathcal A_{D\cup D'}$ and
    their entries vanish. Otherwise the single terms of $(AB)_{\sigma\tau}$ and
    $(BA)_{\sigma\tau}$ are products of equal entries by
    Definition~\ref{def:sites_supported}.
\end{proof}

\section{Placement of a regional matrix}
\label{sec:sites_place}

\begin{definition}[Placement]
    \label{def:sites_place}
    \lean{ManyBody.place, ManyBody.place_apply, ManyBody.reindex_cutEquiv_place,
        ManyBody.place_apply_cutEquiv_symm}
    \leanok
    \uses{def:sites_config}
    The placement of a matrix $K$ on $\C^{X_D}$ is the operator $K\otimes1$ on
    $\mathcal H_X=\C^{X_D}\otimes\C^{X_{V\setminus D}}$, read in the
    coordinates of $\mathcal H_X$ through the cut along $D$. Its matrix element
    $(K\otimes1)_{\sigma\tau}$ is $K_{\sigma|_D,\tau|_D}$ if $\sigma$ and $\tau$
    agree outside $D$, and $0$ otherwise.
\end{definition}

\begin{lemma}[Placement is a unital $*$-homomorphism]
    \label{lem:sites_place_hom}
    \lean{ManyBody.placeHom, ManyBody.placeHom_apply, ManyBody.place_zero,
        ManyBody.place_add, ManyBody.place_smul, ManyBody.place_neg,
        ManyBody.place_sub, ManyBody.place_sum, ManyBody.place_conjTranspose,
        ManyBody.isHermitian_place, ManyBody.place_one, ManyBody.place_mul,
        ManyBody.place_pow, ManyBody.place_mul_place_inv,
        ManyBody.place_inv_mul_place, ManyBody.trace_place}
    \leanok
    \uses{def:sites_place}
    The map $K\mapsto K\otimes1$ is linear, unital, multiplicative and
    compatible with adjoints. If $K$ is invertible, then $K^{-1}\otimes1$ is
    the inverse of $K\otimes1$, and $\Tr(K\otimes1)=|X_{V\setminus D}|\Tr K$.
\end{lemma}
\begin{proof}\leanok
    The cut is a bijection of index sets, so reindexing along it is a
    $*$-isomorphism of matrix algebras that preserves the trace; the
    statements hold for $K\otimes1$ on $\C^{X_D}\otimes\C^{X_{V\setminus D}}$.
\end{proof}

\begin{theorem}[Supported operators are the placed matrices]
    \label{thm:sites_place_range}
    \lean{ManyBody.place_mem_supportedOperators,
        ManyBody.exists_eq_place_of_mem_supportedOperators,
        ManyBody.mem_supportedOperators_iff_exists_eq_place,
        ManyBody.range_placeHom}
    \leanok
    \uses{def:sites_config, def:sites_supported, thm:sites_supported_subalgebra,
        def:sites_place, lem:sites_place_hom}
    An operator on $\mathcal H_X$ is supported on $D$ if and only if it is
    $K\otimes1$ for a matrix $K$ on $\C^{X_D}$. Hence $\mathcal A_D$ is the
    image of the placement.
\end{theorem}
\begin{proof}\leanok
    The matrix elements of $K\otimes1$ given in
    Definition~\ref{def:sites_place} satisfy both conditions of
    Definition~\ref{def:sites_supported}. Conversely, if $\mathcal H_X\neq0$, fix a
    configuration $\sigma_0$ and put $K_{ab}=A_{(a,\sigma_0|_{V\setminus D}),
    (b,\sigma_0|_{V\setminus D})}$; the support conditions give
    $A=K\otimes1$ entrywise. If $\mathcal H_X=0$, then $A=0=0\otimes1$.
\end{proof}

\begin{theorem}[Supported operators are spanned by product operators]
    \label{thm:sites_product_span}
    \lean{ManyBody.supportedOperators_toSubmodule_eq_span,
        ManyBody.prod_mem_supportedOperators, ManyBody.place_single}
    \leanok
    \uses{thm:sites_place_range, lem:sites_place_hom}
    The algebra $\mathcal A_D$ is the linear span of the product operators
    $\bigotimes_{v\in V}m_v$ with $m_v=1$ for every $v\notin D$.
\end{theorem}
\begin{proof}\leanok
    The matrix element of $\bigotimes_vm_v$ at $(\sigma,\tau)$ is
    $\prod_v(m_v)_{\sigma_v\tau_v}$, and the factors at $v\notin D$ are
    $\delta_{\sigma_v\tau_v}$, so these operators are supported on $D$.
    Conversely, by Theorem~\ref{thm:sites_place_range} every supported operator
    is $K\otimes1=\sum_{a,b\in X_D}K_{ab}\,(E_{ab}\otimes1)$, and the placed
    matrix unit $E_{ab}\otimes1$ is the product operator with
    $m_v=E_{a_vb_v}$ for $v\in D$ and $m_v=1$ for $v\notin D$.
\end{proof}

\begin{lemma}[Norm and order under placement]
    \label{lem:sites_place_norm_order}
    \lean{ManyBody.norm_place_le, ManyBody.posSemidef_place,
        ManyBody.place_le_place, ManyBody.place_injective, ManyBody.place_inj,
        ManyBody.norm_place, ManyBody.posSemidef_place_iff,
        ManyBody.place_le_place_iff}
    \leanok
    \uses{def:sites_place, lem:sites_place_hom}
    The placement does not increase the operator norm and preserves positive
    semidefiniteness and the order. If every $X_v$ is nonempty, then the
    placement is injective, $\lVert K\otimes1\rVert=\lVert K\rVert$, and
    $K\otimes1\geq0$ if and only if $K\geq0$, and likewise for the order.
\end{lemma}
\begin{proof}\leanok
    A $*$-homomorphism of finite-dimensional C*-algebras is contractive, and an
    injective one is isometric. The tensor product of positive semidefinite
    matrices is positive semidefinite. If every $X_v$ is nonempty, fix a
    configuration $y$ of $V\setminus D$: the compression of $K\otimes1$ to the
    configurations $(a,y)$, $a\in X_D$, is $K$, which gives injectivity and
    the converse for positivity.
\end{proof}

\section{The expectation onto a region}
\label{sec:sites_region_expectation}

\begin{definition}[Region expectation]
    \label{def:sites_region_expectation}
    \lean{ManyBody.regionExpectation, ManyBody.regionExpectation_apply,
        ManyBody.regionExpectationLM, ManyBody.regionExpectationLM_apply}
    \leanok
    \uses{def:sites_place}
    For a region $D$ put $c=|X_{V\setminus D}|$. The region expectation of an
    operator $A$ on $\mathcal H_X$ is
    \begin{align}
        E_D(A)=c^{-1}\,\Tr_{V\setminus D}(A)\otimes1,
        \label{eq:sites_region_expectation}
    \end{align}
    the normalized partial trace over the sites outside $D$, placed on $D$. Its
    matrix element at $(\sigma,\tau)$ is $0$ unless $\sigma$ and $\tau$ agree
    outside $D$, and is then $c^{-1}\sum_{y}A_{(\sigma|_D,y),(\tau|_D,y)}$, the
    sum running over the configurations $y$ of $V\setminus D$. This is the expectation onto the operators supported on a
    ball in the quasi-local approximation of the two-dimensional area law, there
    for one local dimension.
\end{definition}

\begin{theorem}[The region expectation is a conditional expectation]
    \label{thm:sites_region_expectation}
    \lean{ManyBody.regionExpectation_mem_supportedOperators,
        ManyBody.regionExpectation_mul_mul, ManyBody.regionExpectation_one,
        ManyBody.regionExpectation_of_mem_supportedOperators,
        ManyBody.regionExpectation_regionExpectation,
        ManyBody.trace_regionExpectation, ManyBody.regionExpectation_conjTranspose,
        ManyBody.posSemidef_regionExpectation, ManyBody.regionExpectation_mono,
        ManyBody.norm_regionExpectation_le, ManyBody.regionExpectation_zero,
        ManyBody.regionExpectation_add, ManyBody.regionExpectation_smul,
        ManyBody.regionExpectation_sub, ManyBody.isHermitian_regionExpectation}
    \leanok
    \uses{def:sites_region_expectation, thm:sites_place_range, lem:sites_place_hom,
        thm:root_channel}
    The values of $E_D$ lie in $\mathcal A_D$. The map $E_D$ is linear and
    compatible with adjoints, so it maps Hermitian operators to Hermitian operators. It satisfies
    $E_D(BAC)=B\,E_D(A)\,C$ for all $B,C\in\mathcal A_D$ and $E_D(1)=1$, so it
    fixes $\mathcal A_D$ and $E_D\circ E_D=E_D$. It preserves the trace, maps
    positive semidefinite operators to positive semidefinite operators, is
    monotone, and satisfies $\lVert E_D(A)\rVert\leq\lVert A\rVert$.
\end{theorem}
\begin{proof}\leanok
    By~(\ref{eq:sites_region_expectation}), $E_D(A)$ is a multiple of a placed
    matrix, so it lies in $\mathcal A_D$ by Theorem~\ref{thm:sites_place_range}.
    In the coordinates of the cut along $D$, write $W_{yy'}=c^{-1/2}\,
    (1\otimes\ket{y}\!\bra{y'})$ for configurations $y,y'$ of $V\setminus D$.
    Entrywise,
    \begin{align}
        \sum_{y,y'}(1\otimes\ket{y'}\!\bra{y})\,M\,(1\otimes\ket{y}\!\bra{y'})
            =\Tr_{V\setminus D}(M)\otimes1,
        \label{eq:sites_region_conjugation}
    \end{align}
    so $E_D(A)=\sum_{y,y'}W_{yy'}^*AW_{yy'}$. Each $W_{yy'}$ and its adjoint act
    as the identity on $D$, hence commute with every $K\otimes1$, that is, by
    Theorem~\ref{thm:sites_place_range}, with every element of $\mathcal A_D$;
    this gives $E_D(BAC)=B\,E_D(A)\,C$. From
    Definition~\ref{def:sites_region_expectation},
    $\Tr_{V\setminus D}(1)=c\cdot1$, so $E_D(1)=1$ when $c\neq0$; when $c=0$
    the space $\mathcal H_X$ is zero. Hence $\sum_{y,y'}W_{yy'}^*W_{yy'}=1$, and
    a sum $\sum_pW_p^*AW_p$ with $\sum_pW_p^*W_p=1$ preserves positivity and does
    not increase the operator norm. The trace of the right side
    of~(\ref{eq:sites_region_expectation}) is $c^{-1}\cdot c\Tr A$.
\end{proof}

\begin{theorem}[Localization error through on-site commutators]
    \label{thm:sites_region_expectation_localization}
    \lean{ManyBody.norm_sub_regionExpectation_le}
    \leanok
    \uses{def:sites_region_expectation, def:sites_supported}
    Let $K$ be a region, $B$ an operator on $\mathcal H_X$ and $\varepsilon\colon
    V\to\R$. Suppose that for every site $y\notin K$ and every unitary
    $U\in\mathcal A_{\{y\}}$ we have $\lVert BU-UB\rVert\leq\varepsilon(y)$.
    Then
    \begin{align*}
        \lVert B-E_K(B)\rVert\leq\sum_{y\notin K}\varepsilon(y).
    \end{align*}
\end{theorem}
\begin{proof}\leanok
    \uses{thm:sites_region_expectation, thm:sites_supported_subalgebra}
    If $\mathcal H_X=0$ both sides are trivial, since $\varepsilon(y)\geq
    \lVert B\cdot1-1\cdot B\rVert=0$. Otherwise every $X_v$ is nonempty.
    Enumerate $X_v$ as $\{0,\dots,N_v-1\}$, and for a sign vector
    $s\in\{\pm1\}^{X_v}$ and $k\in\Z/N_v$ let $T_{s,k}$ map the basis vector
    $j$ to $s_i\,i$, where $i=j+k$. Each $T_{s,k}$ is unitary, and
    \begin{align*}
        \sum_{s,k}(T_{s,k})_{ij}\,\overline{(T_{s,k})_{i'j'}}
            =2^{N_v}\,\delta_{ii'}\,\delta_{jj'}:
    \end{align*}
    the sum over the signs vanishes unless $i=i'$, and then exactly one
    translation maps $j$ to $i$. For a family $g=(s_v,k_v)_{v\in V}$ let
    $W_g$ be the product of the operators $T_{s_v,k_v}$ at the sites outside
    $K$ and of the identity on $K$. Since products of product operators are
    formed sitewise, the average of $W_g\,B\,W_g^*$ over all $g$ is, entrywise,
    a constant multiple of $E_K(B)$, and the constant is $1$ by $E_K(1)=1$.
    Hence $B-E_K(B)$ is the average of $(BW_g-W_gB)W_g^*$, whose norm is
    $\lVert BW_g-W_gB\rVert$. Writing $W_g$ as the product of its on-site
    factors, which are unitaries in $\mathcal A_{\{y\}}$, and using
    $B(UW)-(UW)B=(BU-UB)W+U(BW-WB)$ repeatedly bounds this by
    $\sum_{y\notin K}\varepsilon(y)$. This is the argument
    of~\cite[proof of Lemma~4.1]{OpenAI2026AreaLaw}, there for one local
    dimension.
\end{proof}

\begin{corollary}[The commutant of the complement]
    \label{cor:sites_commutant}
    \lean{ManyBody.mem_supportedOperators_iff_forall_commute}
    \leanok
    \uses{def:sites_supported}
    An operator lies in $\mathcal A_D$ if and only if it commutes with every
    element of $\mathcal A_{V\setminus D}$.
\end{corollary}
\begin{proof}\leanok
    \uses{lem:sites_supported_commute, lem:sites_supported_mono,
        thm:sites_region_expectation_localization, thm:sites_region_expectation}
    Operators on disjoint regions commute
    (Lemma~\ref{lem:sites_supported_commute}). Conversely, if $A$ commutes
    with $\mathcal A_{V\setminus D}$, it commutes in particular with every
    unitary supported on a site $y\notin D$, so
    Theorem~\ref{thm:sites_region_expectation_localization} with
    $\varepsilon=0$ gives $A=E_D(A)\in\mathcal A_D$.
\end{proof}

\section{Regional pure states and relabelling}

\begin{definition}[Regional state and entropy]
    \label{def:sites-regional-state}
    \lean{ManyBody.cutVector, ManyBody.regionState, ManyBody.regionEntropy,
        Entropy.regionState_eq_partialTraceRight, Entropy.regionEntropy_eq_vonNeumannEntropy}
    \leanok
    For a vector $\psi$ on a finite system with basis types $X_v$, let
    $\rho_D=\operatorname{Tr}_{D^c}|\psi\rangle\langle\psi|$ and
    $S_\psi(D)=-\operatorname{Tr}(\rho_D\log\rho_D)$, with $0\log0=0$.
    Neither the vector nor its regional state is required to be normalized.
\end{definition}

\begin{lemma}[Regional expectations]
    \label{lem:sites-regional-expectations}
    \lean{Matrix.inner_toEuclideanLin_eq_trace, Matrix.eq_of_forall_trace_mul_eq,
        ManyBody.norm_cutVector, ManyBody.regionState_posSemidef,
        ManyBody.regionState_isHermitian, ManyBody.trace_regionState,
        ManyBody.inner_place, ManyBody.trace_regionState_submatrix_mul,
        ManyBody.inner_place_conj_eq, ManyBody.inner_conj_eq_of_mem_supportedOperators}
    \leanok
    \uses{def:sites-regional-state}
    The matrix $\rho_D$ is positive semidefinite and has trace $\|\psi\|^2$.
    For every matrix $A$ on $D$,
    $\langle\psi,(A\otimes I)\psi\rangle=\operatorname{Tr}(\rho_D A)$.
    If $K$ is invertible and commutes with $\rho_D$, conjugating an observable
    supported on $D$ by $K\otimes I$ preserves its expectation in $\psi$.
\end{lemma}
\begin{proof}
    \leanok
    Positivity and trace preservation follow from the partial trace.
    The expectation identity follows by summing the matrix entries in cut
    coordinates. The final assertion follows by cyclicity of the trace.
\end{proof}

\begin{lemma}[Disjoint regions]
    \label{lem:sites-regional-union}
    \lean{ManyBody.regionUnionEquiv, ManyBody.regionUnionEquiv_apply_fst,
        ManyBody.regionUnionEquiv_apply_snd, ManyBody.place_union_kronecker,
        ManyBody.partialTraceRight_regionState_union}
    \leanok
    \uses{def:sites-regional-state}
    For disjoint regions $B,C$, their union has configuration space
    $X_B\times X_C$. Placing $A\otimes K$ on the union equals the product
    of the placements of $A$ and $K$. The partial trace over $C$ of the
    regional state on $B\cup C$ equals $\rho_B$.
\end{lemma}
\begin{proof}
    \leanok
    The placement identity follows entrywise. Pairing the proposed marginal
    with every matrix on $B$ proves the partial-trace identity.
\end{proof}

\begin{lemma}[Covariance under site and basis bijections]
    \label{lem:sites-regional-relabel}
    \lean{ManyBody.siteConfigurationEquiv, ManyBody.regionConfigurationEquiv,
        ManyBody.complementConfigurationEquiv, ManyBody.relabelVector,
        ManyBody.norm_relabelVector, ManyBody.regionState_relabel,
        ManyBody.regionEntropy_relabel}
    \leanok
    \uses{def:sites-regional-state}
    A bijection of sites, together with bijections of their basis types,
    preserves vector norms. For corresponding regions, the regional matrices
    are related by the induced basis permutation and have equal entropy.
    These statements hold for empty basis types and unnormalized vectors.
\end{lemma}
\begin{proof}
    \leanok
    Reindex the sum over complementary configurations. The resulting regional
    matrices are unitarily conjugate and have the same eigenvalues.
\end{proof}
